%==============================================================================
%  APÉNDICE C — GLOSARIO DE DEFINICIONES LEGALES
%==============================================================================
\chapter{Glosario de definiciones legales}
\label{ap:glosario}

\begin{alerta}[Fuente y vigencia]
Las definiciones de este glosario se transcriben o resumen del \destacar{artículo 2 de la Ley
\Nro 20.720}, en su texto refundido vigente ---última versión de 30 de diciembre de 2023, que
incorpora las reformas de las Leyes \Nro 21.563 y \Nro 21.595---. Cuando el término tiene, además,
un uso doctrinal o forense propio, se añade una nota entre corchetes. Los reenvíos remiten al
capítulo donde el concepto se desarrolla.
\end{alerta}

\begin{description}[leftmargin=0pt,style=nextline,font=\sffamily\bfseries]

\item[Acuerdo de Reorganización Judicial (art. 2 \Nro 1)] El que se suscribe entre una empresa
deudora y sus acreedores para reestructurar sus activos y pasivos, con sujeción a los Títulos 1 y 2
del Capítulo III (y al Título 3 del Capítulo V, para el simplificado). La ley lo llama
indistintamente «Acuerdo». \emph{Véase} capítulo \ref{cap:reorganizacion-ordinaria}.

\item[Acuerdo de Reorganización Extrajudicial o Simplificado (art. 2 \Nro 2)] El que se suscribe
entre una empresa deudora y sus acreedores y se somete a aprobación judicial conforme al Título 3
del Capítulo III. \emph{Véase} capítulo \ref{cap:reorganizacion-simplificada}.

\item[Audiencia Inicial (art. 2 \Nro 4)] La que se celebra ante el tribunal, con presencia del
deudor si comparece, en la liquidación forzosa (\art{120}). \emph{Véase} capítulo
\ref{cap:liquidacion-ordinaria}.

\item[Audiencia de Prueba (art. 2 \Nro 5)] La del juicio de oposición en que se rinden las pruebas
ofrecidas en la Audiencia Inicial (\art{126}).

\item[Audiencia de Fallo (art. 2 \Nro 6)] Aquella en que se notifica la sentencia definitiva que
pone término al juicio de oposición (\art{127}).

\item[Boletín Concursal (art. 2 \Nro 7)] Plataforma electrónica a cargo de la \superir{}, de libre
acceso y gratuita, en que se publican todas las resoluciones y actuaciones de los procedimientos
concursales, salvo que la ley ordene otra forma de notificación. \emph{Es el medio de notificación
por antonomasia del sistema.} \emph{Véase} capítulo \ref{cap:institucionalidad}.

\item[Certificado de Nominación (art. 2 \Nro 8)] El emitido por la \superir{} en que consta la
nominación del veedor o liquidador, titular y suplente.

\item[Comisión de Acreedores (art. 2 \Nro 9)] La que puede designarse para supervigilar el
cumplimiento del acuerdo (en reorganización) o para adoptar los acuerdos que la junta le delegue (en
liquidación, \art{202}).

\item[Cuenta Final de Administración (art. 2 \Nro 11)] La rendición de cuentas que veedor y
liquidador deben efectuar ante el tribunal, observando la normativa contable, tributaria y
financiera. \emph{Es la última ventana de control de la responsabilidad del auxiliar} (capítulo
\ref{cap:institucionalidad}).

\item[Deudor (art. 2 \Nro 12)] Toda empresa deudora o persona deudora, según el procedimiento y la
disposición de que se trate.

\item[Empresa Deudora (art. 2 \Nro 13)] Toda persona jurídica de derecho privado, con o sin fines de
lucro, y toda persona natural que \destacar{dentro de los veinticuatro meses anteriores} al inicio
del procedimiento haya sido contribuyente de primera categoría. \emph{La calificación fija la vía
aplicable} (capítulo \ref{cap:estatuto-deudores}).

\item[Junta de Acreedores (art. 2 \Nro 15)] Órgano concursal integrado por los acreedores del
deudor; puede ser Constitutiva, Ordinaria o Extraordinaria. \emph{Véase} capítulo
\ref{cap:liquidacion-ordinaria}.

\item[Ley (art. 2 \Nro 16)] La Ley de Reorganización y Liquidación de Activos de Empresas y
Personas ---la propia Ley \Nro 20.720---.

\item[Liquidación Forzosa (art. 2 \Nro 17)] La demanda de liquidación presentada por un acreedor,
conforme al Capítulo IV (empresa) o al Capítulo V (persona). \emph{Véase} capítulo
\ref{cap:liquidacion-ordinaria}.

\item[Liquidador (art. 2 \Nro 19)] Persona natural fiscalizada por la \superir{} cuya misión
principal es realizar el activo del deudor y propender al pago de los créditos.

\item[Martillero Concursal (art. 2 \Nro 20)] Martillero público que voluntariamente se somete a la
fiscalización de la \superir{} para realizar los bienes del deudor.

\item[Nóminas (art. 2, \Nro 21 a 24)] Registros públicos que lleva la \superir{} de veedores,
liquidadores, árbitros concursales y martilleros concursales.

\item[Persona Deudora (art. 2 \Nro 25)] Toda persona natural \destacar{no comprendida} en la
definición de empresa deudora ---definición residual---. \emph{Es el sujeto del 93\,\% de los
procedimientos} (capítulo \ref{cap:renegociacion}).

\item[Persona Relacionada (art. 2 \Nro 26)] El cónyuge, los ascendientes, descendientes y
colaterales por consanguinidad o afinidad \destacar{hasta el sexto grado} inclusive y las sociedades
en que participen (salvo las inscritas en el Registro de Valores), y quienes se encuentren en las
situaciones del artículo 100 de la \leyn{18.045}{de Mercado de Valores}. \emph{Su calificación
gatilla posposición, exclusión de quórum y plazo revocatorio agravado} (capítulo \ref{cap:grupos}).

\item[Procedimiento Concursal (art. 2 \Nro 27)] Cualquiera de los regulados en la ley:
Reorganización y Liquidación de la empresa deudora; Renegociación de la persona deudora; y ---tras
la reforma--- Liquidación Simplificada (\Nro 28 A) y Reorganización Simplificada (\Nro 29 A).

\item[Protección Financiera Concursal (art. 2 \Nro 31)] Período que la ley otorga al deudor sometido
a reorganización (ordinaria o simplificada), durante el cual no puede solicitarse ni declararse su
liquidación, ni iniciarse juicios ejecutivos, ejecuciones o restituciones en arrendamiento.
Comprende entre la notificación de la Resolución de Reorganización y el acuerdo, o el plazo legal si
no se alcanza. \emph{Véase} capítulo \ref{cap:reorganizacion-ordinaria}.

\item[Quórum Especial (art. 2 \Nro 32)] \destacar{Dos tercios} del pasivo total con derecho a voto
verificado o reconocido.

\item[Quórum Calificado (art. 2 \Nro 33)] \destacar{Mayoría absoluta} del pasivo total con derecho a
voto verificado o reconocido.

\item[Quórum Simple (art. 2 \Nro 34)] \destacar{Mayoría} del pasivo con derecho a voto
\emph{presente} en la junta. \emph{La distinción entre los tres quórums es decisiva para el cómputo
de cada acuerdo.}

\item[Resolución de Admisibilidad (art. 2 \Nro 35)] Resolución \emph{administrativa} de la \superir{}
en la renegociación (\art{263}), que produce los efectos del \art{264}.

\item[Resolución de Liquidación (art. 2 \Nro 36)] Resolución \emph{judicial} que produce los efectos
del Párrafo 4 del Título 1 del Capítulo IV ---señaladamente el desasimiento (\art{130})---.

\item[Resolución de Reorganización (art. 2 \Nro 37)] Resolución judicial que produce los efectos del
\art{57} (o del \art{286 B} en el simplificado), en particular la protección financiera concursal.

\item[Servicios de Utilidad Pública (art. 2 \Nro 38)] Los regulados por leyes especiales y sujetos a
fiscalización ---agua, electricidad, gas, teléfono, internet---, con régimen propio de continuidad
(\art{171}).

\item[Superintendencia (art. 2 \Nro 39)] La Superintendencia de Insolvencia y Reemprendimiento.

\item[Veedor (art. 2 \Nro 40)] Persona natural fiscalizada por la \superir{} cuya misión principal es
propiciar los acuerdos entre el deudor y sus acreedores, facilitar la propuesta de acuerdo y
resguardar los intereses de los acreedores.

\end{description}

\begin{foco}[Empresa deudora y persona deudora: la definición que ordena todo el sistema]
La bisagra de la Ley \Nro 20.720 está en los \Nro 13 y 25 del artículo 2. \emph{Empresa deudora}
es toda persona jurídica privada y toda persona natural contribuyente de primera categoría en los
últimos veinticuatro meses; \emph{persona deudora} es, por descarte, toda otra persona natural. De
esa dicotomía dependen la vía procedimental, las exigencias de admisibilidad y el régimen recursivo.
Calificar mal al deudor ---y presentar la vía equivocada--- expone la solicitud a una declaración de
incompetencia (capítulo \ref{cap:estatuto-deudores}).
\end{foco}
